\section{Discussion}
\label{sec:discussion}

\subsection{What a monthly strip identifies}

The results support a narrow reading of what calibration to a baseload strip
accomplishes. It fixes six linear functionals of a curve with several thousand
degrees of freedom, and through Proposition~\ref{prop:centering} it fixes the
first moment of the price distribution at every horizon. It does not fix the
intraday shape of the curve, the level over unquoted months, the volatility of
deviations around the curve, or any risk premium. The sensitivity ordering in
Table~\ref{tab:sensitivity} makes the consequence concrete: the largest single
driver of the reported option value is an input that no observed price
constrains.

A reporting convention follows. Rather than a point value, valuation in markets
of this type is better expressed as a value conditional on an explicitly stated
anchor, accompanied by the range induced by plausible alternatives. The $\pm20\%$
anchor band used here produces a call range of $17.80$ to $528.96$ \TRY{} around
a central $166.75$ \TRY, and that interval carries more information than its
midpoint. Section~\ref{sec:multidate} supplies the band with a warrant it
previously lacked: across seven valuation dates the realised anchor error has a
standard deviation of $11.7\%$, close to the $11.5\%$ a uniform $\pm20\%$ prior
implies, though displaced by $-3.5\%$ from zero and not containing the tail.

\subsection{What the strip predicts, and what that costs}

The multi-date evidence bears on a question wider than valuation mechanics.
Whether forward prices are unbiased predictors of subsequent spot outcomes is an
old subject in electricity markets, where limited storage and inelastic demand
make the hedging-pressure equilibrium of \citet{BessembinderLemmon2002} a natural
frame and where realised premia have been measured for mature markets by
\citet{LongstaffWang2004} and \citet{Algieri2021}. On Turkish monthly baseload
contracts over 2023--2026, the answer is that the strip is neither unbiased nor
reliably biased in one direction. Its pooled error averages $-18\%$ of the
forward level, so a buyer of the strip paid on average well above what delivery
subsequently cost; but two of seven dates reversed the sign, and the date-level
mean is not distinguishable from zero at conventional levels once the dependence
between overlapping windows is respected.

For anyone using such a curve, the practical consequence is about capital rather
than about point forecasts. A dispersion of $775$ \TRY{} around a mean of
$-603$ \TRY{} on a price level near $2{,}500$ \TRY{} means the position risk of a
forward-anchored book is dominated by the level of the curve it is anchored to,
not by the shape of the distribution placed around that level. Provisioning
against the residual volatility of Section~\ref{sec:results}, terminal dispersion
of roughly $530$ \TRY{} at the horizons studied, while treating the anchor level
as known would understate the exposure by a wide margin. The ordering in
Table~\ref{tab:sensitivity} says the same thing from the pricing side, and the
multi-date distribution is what puts a number on it.

We stop short of decomposing the error into a risk premium and a forecast error.
Separating the two requires either an instrument for expectations or a structural
model of the supply stack, and seven valuation dates would not identify it. That
the strip was systematically above realised outcomes in the years when hydrology
was favourable, and below in the years when it was not, is consistent with both a
premium and a persistent forecasting difficulty, and this sample cannot
distinguish them.

\subsection{Separation of concerns as a design property}

The centring identity does more than tidy the algebra. It creates a firewall
between the market-identified and the assumed parts of the model, and the value
of that firewall emerged in a way that had not been anticipated when it was
adopted. When the inherited mean-reversion rate was found to be wrong by four
orders of magnitude, the correction changed option values by a factor of four
while leaving the forward reproduction, the monthly fit errors and the realised
level bias unchanged. In a specification where residual dynamics fed back into
the implied forward, everything would have required recalibration and the
diagnosis would have been considerably harder, since the symptom would have
surfaced in the calibration layer rather than only in the dispersion.

The same property holds for the measure change. Because the drift perturbation
enters both the moment system and the pricing equation through a shared field,
$\E^{\Q}[P_t] = F(t)$ survives any $(a_0,a_1)$, verified numerically at
magnitudes far outside any economically plausible range. Risk premium
assumptions can therefore be explored without disturbing the calibration, which
is what makes them reportable as sensitivities rather than as a separate model.

\subsection{Persistence, transformation and nominal prices}

The mean-reversion episode carries a methodological lesson that generalises
beyond this dataset. Reported persistence statistics from a regime-switching
estimation describe several distinct objects: within-regime autoregressive
persistence, the marginal persistence of the observed series, and the persistence
of deviations around a reference curve. These coincide only under restrictive
conditions. Where the estimation sample also spans substantial currency
depreciation, the raw series carries a monetary trend that inflates measured
persistence further, and a coefficient transferred from that context into a
residual-dynamics role will overstate memory by whatever the trend contributes.
Diagnosis here required reconstructing the marginal autocorrelation of the
switching process analytically and verifying it by simulation, since neither
reported statistic alone identified the problem.

This is likely a common failure mode wherever parameters are handed between
modelling stages, and particularly so in emerging markets where nominal series
are non-stationary for reasons unrelated to the physics of the power system. The
defence that eventually worked is worth stating as a rule: every transferred
parameter should name the process it was estimated on, and any mismatch between
that process and its intended use should be treated as an error until shown
otherwise.
